\documentclass[11pt,a4paper]{article}
\usepackage[T1]{fontenc}
\usepackage{lmodern,amsmath,amssymb,amsthm,mathtools,booktabs,microtype}
\usepackage[margin=26mm]{geometry}
\usepackage[colorlinks=true,linkcolor=blue,citecolor=blue,urlcolor=blue]{hyperref}
\setlength{\emergencystretch}{2em}
\newtheorem{theorem}{Theorem}[section]
\newtheorem{lemma}[theorem]{Lemma}
\newtheorem{proposition}[theorem]{Proposition}
\newtheorem{corollary}[theorem]{Corollary}
\theoremstyle{remark}\newtheorem{remark}[theorem]{Remark}
\newcommand{\Q}{\mathbb Q}\newcommand{\Z}{\mathbb Z}
\DeclareMathOperator{\Tr}{Tr}\DeclareMathOperator{\Nm}{N}
\title{A quartic in 76 variables with separated detector valuations}
\author{}\date{16 November 2026}
\begin{document}\maketitle
\begin{abstract}
Under ordinary rational $abc$ and the specified elliptic-group,
simultaneous-prime, trace-height and pole-realization inputs, we construct
a defining quartic for the integers in 76 total rational variables.
The four terms of each detector factor have distinct valuation classes
modulo four. A square rescaling of the pole input then removes one
weight coordinate from each detector. Contact-zero strata are retained
and handled by bounded height, so no nonzero witness is added.
The polynomial is a sum of 66 primitive integral quadratic squares,
has 1250 monomials, and has an exact coefficient-height formula.
Its sufficient fixed coefficient shift remains ineffective.
\end{abstract}

\section{Statement and arithmetic inputs}
Every variable count includes the input. For a reduced rational number
$u/v$, with $v>0$, put $h(u/v)=\log\max(|u|,v)$.
Let
\[
 K=\Q(\rho),\quad\rho^2=2,\quad D=11-6\rho,
 \quad E:y^2=x^3-Dx,
\]
\[
 R_0=(1+\rho,2\rho),\qquad
 P=2R_0=(a_P,b_P)
 =\left(\frac{57}{8}-\frac72\rho,
          -\frac{165}{8}+\frac{417}{32}\rho\right).
\]
We retain the exact arithmetic inputs of~\cite{ABC78}: the equality
$2E(K)=\Z P$ with $P$ non-torsion; the fixed-data simultaneous-prime
theorem with its prescribed local orders, finite avoidance, square
denominator and real approximation; the all-point trace-height theorem;
and the prepared simultaneous congruent-number pole realization
of~\cite{Pole}. Ordinary rational $abc$ is additionally assumed.
The prime and pole inputs have precisely their scope in~\cite{ABC78};
no new realization or prime-distribution hypothesis is introduced.
The algebraic charts and their exceptional loci are covered by
\cite{Reference,Pair,Trace}.

\begin{theorem}\label{thm:main}
Under these inputs there is a computable family
\[
 Q_B\in\Z[X,Y_1,\ldots,Y_{75}],\qquad
 B\ge81\text{ a fixed odd integer},
\]
such that every member is a sum of 66 primitive quadratic squares,
has total degree exactly four, integer content one and 1250 nonzero
coordinate monomials. Its coefficient height is
\begin{equation}\label{eq:coefficientheight}
 H(Q_B)=\max\{20828275368498373068800,
                 B^2+275933278576697172\}.
\end{equation}
There is a fixed $B_*$ such that every odd $B\ge B_*$ satisfies
\[
 x\in\Z\quad\Longleftrightarrow\quad
 \exists\boldsymbol y\in\Q^{75}:Q_B(x,\boldsymbol y)=0.
\]
The parameter $B$ is a coefficient, not a quantified coordinate.
No numerical sufficient $B_*$, or equality with another family's
sufficient threshold, is asserted.
\end{theorem}

The finite coefficient specification is the full quadratic and quartic
table in \nolinkurl{computation/quartic/abc76_v44}. The eight optimized
reference rows are bound to the balanced source of~\cite{Basis}.
The value $B=81$ specifies an algebraic test member; it is not asserted
to be a defining member.

\section{A modulo-four local detector}
For a discrete valuation $v$ with value group $\Z$, put $v(0)=+\infty$.

\begin{lemma}\label{lem:separation}
Suppose that $v(w)=\nu$ is positive and odd, and that every finite
valuation of $f_0,f_1$ is even. For
\begin{equation}\label{eq:factor}
 \mathcal D=1+\frac{C^4}{w^3}
                  +\frac{f_0^2}{w^2}+\frac{f_1^2}{w},
\end{equation}
the finite summand valuations are distinct modulo four. Thus
$\mathcal D\ne0$ and $v(\mathcal D)$ is their minimum, at most zero.
If $v(\mathcal D)>-\nu/2$, then
\begin{equation}\label{eq:bounds}
 v(C)>\frac58\nu,\qquad
 v(f_0)>\frac34\nu,\qquad
 v(f_1)>\frac14\nu.
\end{equation}
Zero values of $C$ or $f_1$ are permitted.
\end{lemma}
\begin{proof}
The finite orders are
\[
 0,\quad4v(C)-3\nu,\quad2v(f_0)-2\nu,\quad2v(f_1)-\nu.
\]
Their classes are $0,-3\nu,-2\nu,-\nu$ modulo four, in some order
the four distinct classes. A unique term has least valuation. Each
term therefore exceeds $-\nu/2$ when their minimum does, and solving
the three inequalities gives~\eqref{eq:bounds}.
\end{proof}

At a good odd place of $E$, every finite valuation of $1/x(Q)$ is
even. If $v(x)<0$, the curve equation has leading term $x^3$, so
$3v(x)=2v(y)$. If $v(x)>0$, its leading term is $-Dx$, so
$v(x)=2v(y)$. The case $v(x)=0$ is immediate. The reciprocal
coordinate of $O$ is zero. Thus Lemma~\ref{lem:separation} applies
to all points allowed by the source, including half-integral indices
and identity numerator points.

\section{The quadratic source}
A field coordinate means two rational coordinates in the basis
$1,\rho$; a bar denotes conjugation. Retain
\[
\begin{array}{ll}
 \text{input and argument:}&X,\lambda\in\Q;\\
 \text{reference:}&f_H,S,T_H,F_H,G_H\in K;\\
 \text{point pair:}&f,h,g,k,s,F_U,F_V\in K;\\
 \text{ratios and comparison squares:}&L,L_0,C_2,I_2\in K;\\
 \text{height and input powers:}&T,z_2,z_4\in\Q;\\
 \text{second argument:}&t_1\in\Q.
\end{array}
\]
The complete reference equations recover precisely
\[
 H\in E(K)\setminus\{O,(0,0),P,P+(0,0)\},
\]
with $S=y(H)/x(H)$, $T_H=y(H-P)/x(H-P)$,
$f_H=1/x(H)$, $F_H=f_H^2$ and $G_H=1/x(H-P)^2$.
Use their eight primitive rational rows in the exact balanced basis
of~\cite{Basis}. This basis is invertibly equivalent over $\Q$ to
the complete reference system, including its exceptional cases.

Retain the four paired field quadrics
\begin{equation}\label{eq:pair}
 \begin{gathered}
 F_U=f^2,\qquad F_V=g^2,\\
 s(f-g)=h+k,\qquad
 s(h-k)=1-D(F_U+fg+F_V),
 \end{gathered}
\end{equation}
and the four field ratio quadrics
\begin{equation}\label{eq:ratios}
 \begin{aligned}
 hL+fS&=0,&(1-DF_U)L+hS&=0,\\
 kL_0+gT_H&=0,&(1-DF_V)L_0+kT_H&=0.
 \end{aligned}
\end{equation}
The reference values $S,T_H$ cannot vanish simultaneously, since
otherwise both reference points would be two-torsion and so would $P$.
Consequently~\cite[Theorem 1.1]{Pair} shows that these rows imply
both reciprocal curve equations. The full paired-chart theorem
of~\cite{Trace} then gives $U,V\in E(K)$, with identity numerator
points allowed, $U\ne V$, and a finite difference $A=U-V$ satisfying
\[
                         x(A)=s^2+D(f+g).
\]
No point-curve row or nonzero ordinate witness is added.

As affine expressions put
\[
 \beta=z_2+3,\qquad C=L-\overline{L_0},\qquad I=L-\beta.
\]
Impose
\begin{equation}\label{eq:commonrows}
 \begin{gathered}
 C_2=C^2,\quad I_2=I^2,\qquad
 T=\Tr\bigl(s^2+D(f+g)\bigr),\\
 z_2=X^2,\quad z_4=z_2^2,\qquad
 t_1-\lambda=B+z_4^2+T^2.
 \end{gathered}
\end{equation}
Thus $\beta=X^2+3$, $T=\Tr x(A)$ and the additive scale is
\begin{equation}\label{eq:scale}
                         a=B+X^8+T^2.
\end{equation}

\subsection{The two detector blocks}
For each argument $t=t_0=\lambda$ and $t=t_1$ retain thirteen rational
coordinates
\[
 R,Z,w,w_2,\eta,\widetilde b,e,u,p,v,q,p_2,q_2
\]
and three field coordinates $A_1,A_2,\Omega=\omega_0+\rho\omega_1$.
Impose
\begin{align}
 R&=t^2-2t,\qquad RZ+Z-R^2+35R=0,\notag\\
 (R+1)w&=Z^2+24Z,\notag\\
 w_2&=w^2,\label{eq:contactrows}\\
 w_2A_1&=ww_2+C_2^2+wF_H+w_2F_U,\notag\\
 w_2A_2&=ww_2+I_2^2+w\overline{G_H}+w_2\overline{F_V},\label{eq:factorrows}\\
 (R+1)\Omega&=A_1A_2,\qquad
 \omega_0^2-2\omega_0\eta-2\omega_1^2=0,\qquad
 w\widetilde b=\eta^2,\label{eq:tracerows}\\
 p_2&=p^2,\qquad e^2+\widetilde b p_2=1,\qquad
 uv=e(w_2-\widetilde b),\notag\\
 q_2&=q^2,\qquad u^2-v^2=\widetilde b q_2.\label{eq:polerows}
\end{align}
Every row is quadratic. The two blocks have independent pole witnesses.
Each block has 19 rational coordinates and 17 rational equations.
The total coordinate count is
\[
                 2+10+14+8+3+1+2\cdot19=76.
\]
There are 28 common reference, pair, ratio and comparison equations,
three height/power equations, one argument equation and 34 detector
equations, for a total of 66.

\subsection{The full pole equivalence away from contact zeros}
For $w\ne0$, define
\begin{equation}\label{eq:polemap}
 b=\widetilde b/w^2,\qquad
 p_{\rm old}=wp,\quad u_{\rm old}=u/w,\quad
 v_{\rm old}=v/w,\quad q_{\rm old}=q.
\end{equation}
Then~\eqref{eq:polerows} is exactly the original pole system
\begin{equation}\label{eq:oldpole}
 \begin{gathered}
 e^2+bp_{\rm old}^2=1,\qquad
 u_{\rm old}v_{\rm old}=e(1-b),\\
 u_{\rm old}^2-v_{\rm old}^2=bq_{\rm old}^2.
 \end{gathered}
\end{equation}
The inverse map is $p=p_{\rm old}/w$, $u=wu_{\rm old}$,
$v=wv_{\rm old}$, $q=q_{\rm old}$, $\widetilde b=w^2b$.
Thus the equivalence holds on every $w\ne0$ stratum, including zero
pole ordinates. Only $w$ is divided by. Prepared pole realization
will be applied to $b$, not to $\widetilde b$.

\section{Local extraction and contact-zero strata}
Put $r=R+1$ and $z=Z+37$. The contact rows give
\[
 r=(t-1)^2,\qquad r(z-r)=36,
\]
so $r\ne0$. With
\[
 \Theta=\{-5,-2,-1,0,2,3,4,7\},\qquad
 F(t)=\prod_{\theta\in\Theta}(t-\theta),
\]
they also give
\begin{equation}\label{eq:weight}
                              r^3w=F(t).
\end{equation}
Indeed $r^2Z(Z+24)=(r-1)(r-4)(r-9)(r-36)$.
Consequently $w=0$ is equivalent to $t\in\Theta$.

For $w\ne0$, equations~\eqref{eq:factorrows} give
$A_i=w\mathcal D_i$, where
\begin{equation}\label{eq:twophysical}
 \begin{split}
 \mathcal D_1&=1+C^4/w^3+f_H^2/w^2+f^2/w,\\
 \mathcal D_2&=1+I^4/w^3+\overline{g_H}^{\,2}/w^2
                                     +\overline g^{\,2}/w,
 \end{split}
\end{equation}
and $g_H=1/x(H-P)$. Put $W=\mathcal D_1\mathcal D_2$.
Whenever $W\ne0$, the trace row forces $\Tr\Omega\ne0$:
otherwise it forces $\Nm(\Omega)=0$ and hence $\Omega=0$,
contradicting $r\Omega=w^2W$. Therefore
\begin{equation}\label{eq:physicalpole}
 \tau=\frac{\Nm(W)}{\Tr W},\qquad
 \eta=\frac{w^2}{r}\tau,\qquad
 \widetilde b=\frac{w^3}{r^2}\tau^2,\qquad
 b=\frac{w}{r^2}\tau^2.
\end{equation}

\begin{lemma}\label{lem:local}
At every prime counted by the source contact radical $M(t)$ for a
noncontact argument, $\beta$ is $p$-integral and
$n\equiv\beta\pmod p$, where $n\in\tfrac12\Z$ is the index of
$A=U-V$ relative to $P$.
\end{lemma}
\begin{proof}
Such a prime is outside the fixed exceptional set. It has a positive
odd contact order $\nu=v_p(w)$ and $r$ is a unit. Lemma~\ref{lem:separation}
applies at every place above $p$ to both factors. Thus both are
nonzero and have valuation at most zero; in particular $W\ne0$
and~\eqref{eq:physicalpole} is justified.

The pole obstruction applied through~\eqref{eq:polemap} excludes a
negative odd valuation of $b$. Since
$v_p(b)=\nu+2v_p(\tau)$ is odd, it follows that
$v_p(\tau)>-\nu/2$. The elementary reciprocal-trace inequality
\[
                  v_p(\tau)\le\max_{\mathfrak p\mid p}v_{\mathfrak p}(W)
\]
holds also when trace terms cancel. Choose a place attaining a value
larger than $-\nu/2$. As both factor valuations are at most zero,
each factor separately exceeds that lower bound. Applying
Lemma~\ref{lem:separation} gives positive valuations of both comparisons
and of the four relevant reciprocal abscissas. Thus $H,U$ are formal
at one place and $H-P,V$ at its conjugate place. A zero numerator
reciprocal represents the identity and is included.

Here is the index step, to specify why positive formal depth suffices.
Write the labels of $H,U,V$ as $h_0,u_0,v_0\in\tfrac12\Z$.
The reference labels are $h_0,h_0-1$, at least one a $p$-adic unit.
The formal logarithm extends to the whole local elliptic group by
first multiplying into the formal subgroup and then dividing by that
integer. Its value at a point of label $a$ is $a$ times its nonzero
value at $P$; two-torsion contributes zero. On the formal subgroup
the logarithm has the same valuation as its parameter, and differs
from that parameter by a factor in $1+\mathfrak p$.

At the reference with unit label, the corresponding numerator ratio
is therefore integral, since numerator labels are $p$-integral.
The congruences $C\equiv I\equiv0$ then make $\beta,L$ and
$\overline{L_0}$ integral. If a reference label is a unit, the
logarithm comparison directly yields the corresponding congruence
below. If it is divisible by $p$, integrality of the raw ratio forces
the numerator label to be divisible by $p$ too, giving the same
congruence. Identity numerators give label and ratio zero. Hence
\[
 u_0\equiv h_0L\equiv h_0\beta,\qquad
 v_0\equiv(h_0-1)\overline{L_0}\equiv(h_0-1)\beta\pmod p.
\]
Subtraction gives $n=u_0-v_0\equiv\beta$. The argument applies at
split and inert places and makes no Gaussian prime-selection assumption
about an arbitrary soundness witness.
\end{proof}

At $w=0$ the factor rows instead force $C=I=0$. These strata are
retained; no division by $w$ or local extraction is asserted there.
Their arguments belong to the fixed set $\Theta$, so have height at
most $\log7$. Both arguments cannot be contacts, because their real
difference~\eqref{eq:scale} is at least 81, greater than the contact
diameter 12. The source includes no extra nonzero-coordinate guard.

\section{Height soundness}
We use the proved ordinary-$abc$ consequences in~\cite{ABC78}:
for $t\notin\Theta$,
\begin{equation}\label{eq:contactheight}
 h(t)\le\frac27\log M(t)+C_0,
\end{equation}
and, for $a=B+X^8+T^2$,
\begin{equation}\label{eq:mixedheight}
 h(X)\le\frac{301}{388}h(a)+C_1,\qquad
 h(T)\le\frac{301}{97}h(a)+4C_1,\qquad h(a)\ge\log B.
\end{equation}
All constants are independent of $B,X,T$ and the witnesses.
The source all-point trace theorem gives fixed constants $A_0,J_0$
such that, for $m=2n\in\Z$,
\begin{equation}\label{eq:traceheight}
                         m^2\le A_0h(T)+J_0.
\end{equation}
It includes torsion differences and index zero.

Suppose $n\ne\beta$. Write $X=u_X/v_X$ in lowest terms with
$v_X>0$ and put
\[
 N=|m v_X^2-2u_X^2-6v_X^2|>0.
\]
At a consumed prime, Lemma~\ref{lem:local} gives $p\nmid v_X$
and $p\mid N$. Hence $M(t_j)\mid N$ for every noncontact argument.
For a contact argument use $h(t_j)\le\log7$ instead. Since $N\ge1$,
both cases yield the uniform estimate
\[
 h(t_j)\le\frac27\log N+C_2,
 \qquad h(a)\le\frac47\log N+C_3.
\]
No coprimality of the two contact radicals is needed. Also
\[
                   N\le8(1+|m|)e^{2h(X)}.
\]
Combining these inequalities with~\eqref{eq:mixedheight} gives
\[
 h(X)\le\frac{602}{679}h(X)
             +\frac{301}{679}\log(1+|m|)+C_4.
\]
The strict coefficient is $602/679=86/97<1$. Rearrangement and the
companion trace estimate give
\[
 h(X)\le\frac{301}{77}\log(1+|m|)+C_5,\qquad
 h(T)\le\frac{1204}{77}\log(1+|m|)+C_6.
\]
Equation~\eqref{eq:traceheight} first bounds $|m|$, then $h(X)$,
$\log N$ and $h(a)$, by constants independent of $B$. This contradicts
$h(a)\ge\log B$ for all sufficiently large fixed odd $B$.

Consequently $n=X^2+3$. Thus $2X^2=m-6$ is an integer. In a reduced
fraction $X=u_X/v_X$, this forces $v_X^2\mid2$, so $v_X=1$ and
$X\in\Z$. This proves soundness, including the permitted contact-zero
strata.

\section{Completeness and the order of choices}
Fix $X\in\Z$ and an odd $B\ge81$. Set
\[
 \beta=X^2+3\ge3,\qquad A=[\beta]P,\qquad
 T=\Tr x(A),\qquad a=B+X^8+T^2.
\]
These are fixed before invoking the simultaneous-prime theorem.
The two contact lists $\Theta$ and $\Theta-a$ are disjoint because
$a>12$. Their denominators, all finite preparation data and the chosen
real neighborhoods are fixed before prime selection.

The retained theorem and compatible order congruences supply the argument
$\lambda$ and a common index $h_0$ with
\[
 H=[h_0]P,\qquad U=[h_0\beta]P,\qquad
 V=[(h_0-1)\beta]P,
\]
so $U-V=A$. The source formal estimates give comparison orders at
least four and reciprocal orders at least two at every selected simple
contact. The three nonconstant terms in~\eqref{eq:twophysical}
then have orders at least $13,2,3$ in the intended embedding.
Both factors are consequently units congruent to one.

In the conjugate embedding, each opposite reference reciprocal is a
unit: the reference reductions are $P$ and $-P$ instead of the identity.
The reference term therefore has order $-2$. The modulo-four separation
makes each conjugate factor have negative order, at most $-2$.
Thus $W$ is a unit in the intended embedding and has negative order
in its conjugate. There is no equal-order trace cancellation, and
$\tau$ is a local unit. Formula~\eqref{eq:physicalpole} gives
$v_p(b)=1$ at each selected odd contact.

At each fixed preparation prime, the point and ratio functions range
in fixed bounded local neighborhoods. Increasing the prescribed
denominator depth makes $1/w$ arbitrarily small, so the three new
nonconstant terms tend to zero. Thus $W$ tends to one and
$\Tr(W^{-1})$ tends to two. The same finite preparation as in the
source makes the physical input $b$ a square in $\Q_2$ with
$v_2(b)\ge3$. The odd-valuation primes of $w$ outside the fixed
exceptional set are the selected contacts, where $v_p(b)=1$ as just
proved; all remaining valuations of $b$ are even. Hence its reduced
denominator is a square.

On the fixed real compact neighborhoods, $1/w$ tends to zero as both
arguments increase. The factors tend uniformly to one, and
\[
                 b=\frac{w}{r^2}\tau^2\sim\frac1{4t^2}.
\]
The source real preparation therefore gives $0<b<1$. The retained
simultaneous pole-realization theorem now applies to the two actual
physical inputs $b$, with independent witnesses of~\eqref{eq:oldpole}.
Transport them by~\eqref{eq:polemap} to obtain every row of the new
system. No preparation inequality is required for $\widetilde b$ itself.
The same construction covers $X=0$, for which $\beta=3$.

This proves completeness for every integer and every fixed odd $B\ge81$.
Together with soundness it proves the defining assertion of
Theorem~\ref{thm:main}. The sufficient soundness threshold is ineffective
and is not identified with a threshold for another detector family.

\section{Exact polynomial and verification}
Clear each rational row to a primitive integer quadratic $q_{B,j}$ and set
\[
                         Q_B=\sum_{j=1}^{66}q_{B,j}^2.
\]
Over $\Q$, this has exactly the common zero set of the rows. The input
$X$ occurs only in the primitive row $z_2-X^2$, so its fourth power
has coefficient one in $Q_B$. This proves exact degree four and content
one for every fixed integer $B$.

The producer binds the balanced parent quadratic table by its SHA-256
digest, removes exactly the two cube-weight coordinates and their rows,
derives twelve changed rational rows, and preserves the remaining
54 rows by name. Its full expansion has 1253 parameter-separated terms
and 1250 coordinate monomials. The coefficient profile of the constant
monomial is $B^2+275933278576697172$; the largest fixed coefficient is
$20828275368498373068800$. All other parameter-dependent profiles have
linear coefficient of absolute value at most two and constant part of
absolute value at most one, with no cancellation at positive $B$.
They are bounded by the displayed constant profile. This proves
\eqref{eq:coefficientheight} and the support assertion from the exact
coefficient tables.

The separate checker imports no producer code. It reconstructs the
changed field equations in the basis $1,\rho$, verifies every retained
row and every expanded coefficient with independent symbolic arithmetic,
and regenerates the tables in an isolated directory byte for byte.
It also checks the pole transformation, valuation residue classes,
contact factorization, contact-zero witnesses and strict height coefficient.
These finite certificates verify the polynomial specification and the stated
algebraic identities. The infinite arithmetic conclusions use the proofs
above and the retained hypotheses at their specified scope.

\begin{thebibliography}{9}
\bibitem{ABC78}\emph{A quartic in 78 variables from a mixed-power height
estimate}, project manuscript, 15 November 2026,
\nolinkurl{papers/quartic_existence/quartic_abc78_20261115_v34.tex}.
\bibitem{Reference}\emph{A complete five-coordinate elliptic reference chart},
project manuscript, 14 November 2026,
\nolinkurl{papers/elliptic_charts/compressed_reference_20261114_v27.tex}.
\bibitem{Pair}\emph{Recovery of elliptic curve equations from complete
ratio equations}, project manuscript, 15 November 2026,
\nolinkurl{papers/elliptic_charts/paired_ratio_redundancy_20261115_v40.tex}.
\bibitem{Trace}\emph{Trace heights on quadratic twists with $j=1728$},
project manuscript, 13 November 2026,
\nolinkurl{papers/trace_heights/trace_on_all_points_20261113_v25.tex}.
\bibitem{Basis}\emph{Primitive quadratic bases and small-coefficient quartic
presentations},
project manuscript, 15 November 2026,
\nolinkurl{papers/quartic/primitive_quadratic_bases_20261115_v40.tex}.
\bibitem{Pole}\emph{Pole equations from simultaneous congruent-number
twists}, project manuscript, 14 November 2026,
\nolinkurl{papers/pole_parameters/congruent_twist_poles_20261114_v27.tex}.
\end{thebibliography}
\end{document}
